% from ega4-9-fr.tex
\oldpage[IV-3]{83}

sur l'ouvert $Z\cap X_s$. Alors, dans le théorème (9.7.7), on peut remplacer
$X_s$ par $Z_s$ sans changer la conclusion~: on le voit en répétant le
raisonnement fait dans (9.6.7, (ii)).
\end{remark}

\begin{proposition}[9.7.12]
\label{IV.9.7.12-fr}
Soient $f:X\to S$ un morphisme de présentation finie, et $g:S\to X$ une
$S$-section de $X$ (\textbf{I}, 2.5.5). Pour tout $s\in S$, soit $X_s^0$ la
composante connexe de $f^{-1}(s)$ qui contient $g(s)$~; alors, la réunion $X^0$
des $X_s^0$ pour $s\in S$ est une partie localement constructible de $X$.
\end{proposition}

Montrons d'abord qu'on peut se ramener au cas où $S$ est affine et
\emph{noethérien}. On peut toujours supposer $S=\operatorname{Spec}(A)$ affine~;
avec les notations de (9.2.3), on a $f=(f_0)_{(S)}$, où $f_0:X_0\to S_0$ est un
morphisme de type fini, et on peut en outre supposer qu'il existe une
$S_0$-section $g_0:S_0\to X_0$ telle que $g=(g_0)_{(S)}$ (8.9.1). Remarquons
maintenant que si $p$ est le morphisme $S\to S_0$, alors, pour tout $s_0\in S_0$,
la composante connexe $(X_0)^0_{s_0}$ de $f_0^{-1}(s_0)$ qui contient $g_0(s_0)$
est \emph{géométriquement connexe} (4.5.13), et par suite, si $s_0=p(s)$, on a
$X_s^0=q^{-1}((X_0)^0_{s_0})$ où $q:X_s\to(X_0)_{s_0}$ est la projection
canonique ((4.5.8) et (4.4.1))~; notre assertion résulte donc de (1.8.2).

Utilisons alors le critère de constructibilité $(\mathbf 0_{\mathrm{III}},9.2.3)$~:
soient $x$ un point de $X$, $Z$ le sous-préschéma réduit de $X$ ayant
$\overline{\{x\}}$ comme espace sous-jacent, $Y$ le sous-préschéma réduit de $S$
ayant $\overline{f(Z)}$ comme espace sous-jacent~; comme la restriction de $f$
à $Z$ se factorise en $Z\to Y\to S$ (\textbf{I}, 5.2.2), on peut remplacer $S$
par $Y$, autrement dit supposer que $f(x)=\eta$, point générique du préschéma
intègre $S$. Par hypothèse, $X_\eta$ est somme de deux sous-préschémas
$X_\eta^0$, $X_\eta^1$, induits sur des ouverts complémentaires de $X_\eta$. En
vertu de (9.5.4) et (9.5.1), on peut donc, en remplaçant au besoin $S$ par un
voisinage de $\eta$, supposer que $X$ est réunion de deux ouverts disjoints
$X^0$, $X^1$ tels que $(X^0)_\eta=X_\eta^0$ et $(X^1)_\eta=X_\eta^1$. Comme
$g:S\to X$ est continu et injectif, $S$ est réunion des deux ouverts disjoints
$g^{-1}(X^0)$ et $g^{-1}(X^1)$~; mais comme $S$ est irréductible, et
\emph{a fortiori} connexe, un de ces deux ouverts est vide, et comme
$g(\eta)\in X_\eta^0$ par définition, on a $g^{-1}(X^1)=\varnothing$. En
d'autres termes, $g$ est une $S$-section de $X^0$~; d'autre part, comme
$(X^0)_\eta$ est géométriquement connexe (4.5.13), il en est de même de
$(X^0)_s$ pour tout $s$ voisin de $\eta$ (9.7.7)~; comme
$g(s)\in(X^0)_s$, on a bien $(X^0)_s=X_s^0$. C.Q.F.D.

\subsection{Décomposition primaire au voisinage d'une fibre générique}
\label{subsection:IV.9.8-fr}

\begin{proposition}[9.8.1]
\label{IV.9.8.1-fr}
Soient $f:X\to S$ un morphisme de présentation finie, $\mathcal F$ un
$\mathcal O_X$-Module quasi-cohérent de présentation finie. Alors l'ensemble
$E$ des $s\in S$ tels que $\mathcal F_s$ soit un $\mathcal O_{X_s}$-Module sans
cycle premier associé immergé est localement constructible.
\end{proposition}

On sait que pour les Modules quasi-cohérents sur les préschémas algébriques, la
propriété d'être sans cycle premier associé immergé est invariante par
changement du corps de base (4.2.7)~; on peut donc se borner au cas où $S$ est
affine, noethérien et intègre de point générique $\eta$, et à prouver que dans
ce cas $E$ ou $S-E$ est un voisinage de $\eta$. On a vu dans (9.7.6) que si
$\eta\in E$, $E$ est un voisinage de $\eta$~; reste donc à considérer le cas où
$\mathcal F_\eta$ a des cycles premiers associés immergés. On peut se borner au
cas où $X$ est affine, et où il y a un sous-$\mathcal O_{X_\eta}$-Module
cohérent $\mathcal H$ de $\mathcal F_\eta$ dont le support $Z'$ est non vide et
\emph{rare} par rapport au support $Z$ de $\mathcal F$ (3.1.3)~; alors, en vertu
de (8.5.2, (i) et (ii))

